\section{De la cobertura lingüística a la utilidad real}

Para entender el punto de partida de NLLB, primero hay que distinguir tres cosas que a menudo se confunden: cuántas lenguas escritas existen en el mundo, cuántas lenguas cubrían los productos comerciales en ese momento y cuántas de ellas alcanzaban realmente una calidad de traducción utilizable, confiable y segura. En clase, una serie de slides muy sencillas desarrolla estos tres niveles uno por uno.

\subsection{Tres mil lenguas escritas y una motivación personal}

El primer mapa del mundo no busca un número total de lenguas exacto e invariable, sino establecer un orden de magnitud: las lenguas escritas del mundo superan con mucho lo que cubren los principales productos digitales. También sugiere un problema más profundo: «tener escritura» no equivale a «tener abundante texto digital», y menos aún a contar con corpus paralelos de alta calidad que puedan usarse directamente para entrenar.

\lecturefigure{slide-02-written-languages.jpg}{En el mundo existen más de tres mil lenguas escritas}{Video de la clase de Stanford 00:01:24}

\begin{knowledgebox}{Lectura de la figura: existencia de la lengua, digitalización y aptitud para el entrenamiento son tres condiciones distintas}
Las lenguas del mapa necesitan, en primer lugar, un identificador estable de la lengua o de la variante; en segundo lugar, texto digital accesible; y solo en tercer lugar pueden llegar a tener corpus bilingües alineados a nivel de oración. Una lengua de bajos recursos puede estar muy viva en el primer nivel y ser casi invisible en los otros dos. Por lo tanto, tener pocos datos de entrenamiento no es sinónimo de tener pocos hablantes, de que la lengua valga menos ni de que la lengua sea simple.
\end{knowledgebox}

\begin{warningbox}{No tomar el «número de lenguas escritas» como un valor verdadero fijo}
Los límites entre lengua y dialecto, los sistemas de escritura y las formas de normalización implican decisiones sociales. La cifra de la clase sirve para mostrar la brecha de cobertura; no debe usarse para asignar a la fuerza un código único a cada lengua ni para negar el criterio de las comunidades locales sobre nombres, escrituras y variantes.
\end{warningbox}

\teachervoice{El docente enfatiza que la tecnología multilingüe no es algo que nunca haya existido; la traducción automática es, de hecho, una de las aplicaciones comerciales de la IA generativa más exitosas y extendidas. La cuestión no es «si hay traducción», sino por qué esa capacidad sigue concentrada en unas pocas lenguas.}

\subsection{La cobertura de un producto es una instantánea en el tiempo, no una prueba de calidad}

A continuación, la clase coloca en una misma figura el número de lenguas del mundo y la cobertura que tenían en ese momento Google Translate y Microsoft Translator. El primer propósito de esta comparación es mostrar la diferencia en cantidad; el segundo, recordar que el conteo de lenguas en la página de un producto solo indica que la interfaz las admite, y no dice nada sobre la distribución de errores en textos largos, dominios especializados, direcciones de bajos recursos o escenarios donde la seguridad importa.

\lecturefigure{slide-03-product-coverage.jpg}{Cobertura lingüística de los productos comerciales de traducción en la fecha de la clase}{Video de la clase de Stanford 00:02:03}

\begin{knowledgebox}{Lectura de la figura: primero comparar órdenes de magnitud, después preguntar cómo se define la calidad}
En la figura, la cobertura de los productos es del orden de cien lenguas, mientras que las lenguas escritas son del orden de miles. A primera vista debe notarse la diferencia de cantidad; en una segunda lectura hay que preguntar si cada «lengua admitida» incluye traducción en ambos sentidos, si fue validada por hablantes nativos y si funciona igual de bien en las direcciones de bajos recursos. El número de lenguas es un indicador de coverage, no de quality, safety ni equity.
\end{knowledgebox}

\teachervoice{Al dar el número de lenguas de los productos, Fan aclaró por iniciativa propia que esas cifras «quizá ya estén un poco desactualizadas». Por eso, estas notas conservan como referencia temporal la clase de 2023 y no presentan las cifras de la slide como especificaciones vigentes de los productos en 2026.}

\subsection{El objetivo no es duplicar, sino duplicar de forma segura}

Esta sección convierte la diferencia de cantidad anterior en un objetivo de ingeniería. La pregunta del proyecto NLLB se resume en una frase: si se parte de unas cien lenguas, ¿cómo duplicar la cobertura? Pero el paréntesis que la slide agrega después es el núcleo de la restricción: no basta con añadir etiquetas de lengua; también hay que obtener, en la medida de lo posible, traducciones de alta calidad, seguras y realmente utilizables.

\lecturefigure{slide-04-double-coverage-safely.jpg}{Duplicar la cobertura lingüística manteniendo alta calidad y seguridad}{Video de la clase de Stanford 00:02:51}

\begin{importantbox}{Cómo formular correctamente el objetivo de cobertura}
Sea $\mathcal{L}$ el conjunto de lenguas y $\mathcal{D}\subseteq\mathcal{L}\times\mathcal{L}$ el conjunto de direcciones que el sistema puede afirmar que cubre. El objetivo real no es maximizar solamente $|\mathcal{L}|$ o $|\mathcal{D}|$, sino cumplir en cada dirección restricciones mínimas de calidad, seguridad y utilidad:
\[
\max |\mathcal{D}|\quad
\text{s.t.}\quad q_d\ge \tau_q,\ s_d\ge \tau_s,\ u_d\ge \tau_u\quad \forall d\in\mathcal{D}.
\]
Aquí $q_d$ representa la calidad de la traducción, $s_d$ el desempeño en seguridad y $u_d$ la utilidad para los usuarios; $\tau$ es el umbral mínimo que define el proyecto. La clase no dio umbrales concretos; la fórmula sirve para ilustrar las restricciones de múltiples objetivos, no para inventar reglas de gobernanza.
\end{importantbox}

\begin{warningbox}{La cobertura lingüística no se cuenta por etiquetas}
Una misma lengua puede abarcar distintas escrituras, variantes regionales y dominios; un mismo par de lenguas se divide además en direcciones into-English, out-of-English y non-English directions. Si solo se agregan nombres de lenguas a un menú desplegable, el sistema puede seguir sin ser utilizable justo en las direcciones donde más se necesita.
\end{warningbox}

\subsection{«Altos recursos» describe los datos, no la población}

Al repasar la evolución histórica, la clase señala que la traducción automática se desarrolló durante mucho tiempo en torno a las lenguas de altos recursos, y que corpus paralelos institucionales como Europarl fueron la base inicial. El término high-resource describe principalmente los datos de entrenamiento y evaluación disponibles, no la población que habla la lengua. Una lengua con muchísimos hablantes puede seguir siendo de bajos recursos por falta de texto digital, de normalización o de corpus paralelos.

\lecturefigure{slide-05-current-progress.jpg}{La traducción automática pasa gradualmente de las lenguas de altos recursos a las de bajos recursos}{Video de la clase de Stanford 00:04:00}

\begin{knowledgebox}{Lectura de la figura: el eje de recursos y el eje de población son independientes}
La trayectoria histórica de la izquierda depende de textos bilingües institucionales; los avances recientes de la derecha dependen del modelado multilingüe, la supervisión débil y la minería de corpus. Al leer la figura conviene separar el «volumen de recursos de datos» del «tamaño de la población»: tener muchos recursos suele facilitar el modelado, pero no indica una mayor necesidad social; tener pocos recursos debilita a la vez el entrenamiento, el ajuste de hiperparámetros y la evaluación.
\end{knowledgebox}

\begin{knowledgebox}{Términos clave: recursos y datos}
\begin{tabularx}{\textwidth}{@{}L{0.21\textwidth}L{0.28\textwidth}X@{}}
\toprule
Término & Problema que resuelve & Significado en esta clase \\
\midrule
High-resource language & Describir la suficiencia de datos & Cuenta con abundantes recursos monolingües, bilingües y de evaluación; no equivale a tener más hablantes \\
Low-resource language & Describir la escasez de datos & Faltan recursos de entrenamiento, de normalización y de evaluación humana \\
Monolingual data & Aprender la distribución de una sola lengua & Textos, como páginas web o documentos, que contienen una sola lengua \\
Bitext / parallel data & Aprender la correspondencia entre lenguas & Pares de oraciones de origen y destino alineadas semánticamente \\
Translation direction & Precisar el sentido de entrada y salida & $A\rightarrow B$ y $B\rightarrow A$ son dos tareas distintas \\
\bottomrule
\end{tabularx}
\end{knowledgebox}

\subsection{Resumen de la sección}

El punto de partida de NLLB no es «en el mundo hay tres mil lenguas, así que hay que entrenar tres mil», sino incorporar en el objetivo, de forma conjunta, la cobertura, la calidad, la seguridad y la utilidad. Los altos y bajos recursos describen las condiciones de los datos; el número de lenguas de un producto es solo una instantánea en el tiempo y no sustituye la evaluación por dirección.
